\section{A positive codimension-two bridge through qubit geometry}
Let $V\in\mathbb C^{n\times(n-2)}$ have unit-modulus entries and
$V^*V=nI_{n-2}$. Its residual Gram matrix
\[
 K=nI_n-VV^*
\]
is positive semidefinite of rank two, with diagonal entries two.
It has a factorization $K=RR^*$ with $R\in\mathbb C^{n\times2}$,
$R^*R=nI_2$, and every row of $R$ has squared norm two.
Define the real matrix
\[
 B_{ij}=\frac{|K_{ij}|^2}{2}-1.
\]

\begin{theorem}[An exact rank test for a two-column flat completion]
The matrix $B$ is a Gram matrix of unit vectors in $\mathbb R^3$.
The partial complex Hadamard matrix $V$ can be completed by two
unit-modulus columns if and only if $\operatorname{rank}B\le2$.
\end{theorem}
\begin{proof}
For each row $r_i$ of $R$, let $u_i=r_i^T/\sqrt2$ and form its unit
Bloch vector $b_i\in\mathbb R^3$. The elementary two-dimensional
identity $b_i\cdot b_j=2|u_i^*u_j|^2-1$ gives $B_{ij}$.
All other residual factors are obtained by a common two-dimensional
unitary change of basis. Such a factor has both entry moduli one
in every row exactly when every corresponding qubit has equal
coordinate probabilities, that is, every $b_i$ is perpendicular to
a fixed unit measurement axis. Such an axis exists exactly when
the Bloch vectors span dimension at most two. Explicitly, for a
unit normal $s$, diagonalize the traceless Hermitian matrix
$s\cdot\sigma$, whose eigenvalues are $\pm1$. Zero expectation in
each $u_i$ forces its two eigenbasis probabilities to be $1/2$.
This constructs the required residual factor.
\end{proof}

\begin{theorem}[Quaternary input needs no extra phase freedom]
Suppose $n$ is even and $V$ has entries in $\mu_4$.
Then $B_{ij}\in\{-1,0,1\}$, and $\operatorname{rank}B\le3$.
Whenever $\operatorname{rank}B\le2$, the two completing columns
can also be chosen in $\mu_4$. Thus complex and quaternary
completability are equivalent for this codimension-two input.
\end{theorem}
\begin{proof}
The off-diagonal $K_{ij}$ are Gaussian integers with real and
imaginary parts of even sum, since they are sums of the even number
$n-2$ of fourth roots. Positive semidefiniteness gives $|K_{ij}|\le2$.
Thus their squared magnitudes are $0,2,4$, proving the assertion on
$B$. Distinct Bloch vectors are consequently parallel, antiparallel,
or orthogonal, and lie on at most three mutually orthogonal axes.

If the rank is at most two, rotate their plane to the equator.
After a common phase choice the residual rows have the form
$a_i(1,\delta_i)$ with $|a_i|=1$ and
$\delta_i\in\{1,-1,i,-i\}$. For two nonorthogonal rows,
\[
 K_{ij}=a_i\overline{a_j}(1+\delta_i\overline{\delta_j}).
\]
If the two $\delta$ values agree, the Gaussian integer $K_{ij}$
has modulus two, so division by two is a fourth root. If the axes
differ, its modulus is $\sqrt2$, so division by $1+i$ or $1-i$
is again a fourth root. Therefore $a_i\overline{a_j}\in\mu_4$
on every nonorthogonality edge. For rank two this graph is connected:
states on one axis are nonorthogonal to every state on the other.
A global phase makes all $a_i$ fourth roots. In rank one the two
antipodal groups are disconnected but mutually orthogonal; choose
their common phases independently without changing $K$.
This also makes every $a_i$ a fourth root. The resulting residual
matrix has entries in $\mu_4$ and the same Gram $K$, and appending
it to $V$ gives the required quaternary matrix.
\end{proof}
The proof is an exact algorithm using only $K$: take a reference
Bloch axis from a row of $B$, choose an orthogonal second axis if
needed, and label their antipodes by $1,-1,i,-i$. Along nonzero
$K_{ij}$ edges solve the displayed phase ratio, starting one phase
at one in each connected component. All recovered phases are in
$\mu_4$. This avoids numerical extraction of a Bloch rotation.

\begin{theorem}[Automatic quaternary completion using a classical real input]
Every even-order quaternary partial Hadamard matrix missing two
columns has a quaternary completion.
This assertion uses the classical real codimension-four Hadamard
completion theorem as an input; it is not proved here from scratch.
\end{theorem}
\begin{proof}
The case $n=2$ is immediate. For $n\ge4$, realify $(1+i)V$ to a
real sign matrix of size $2n\times(2n-4)$ with column Gram $2nI$.
Classical real completion supplies four sign columns $Q$ with
\[
 QQ^T=2\mathcal R(K).
\]
On the other hand $\sqrt2\mathcal R(R)$ is a full-column-rank real
factor with that same Gram. The two factors differ by a right
orthogonal matrix, preserving the number of antipodal row directions.

If $B$ had rank three, the previous discrete inner-product argument
would give three orthogonal Bloch axes. Moreover $R^*R=nI_2$
implies $\sum_i b_i=0$, so both antipodes of each axis occur.
There would be at least six distinct complex lines represented by
the rows $r_i$. Realification contributes two different orthogonal
real row directions from each complex line; different complex lines
cannot share a nonzero real direction. Thus
$\sqrt2\mathcal R(R)$, and hence $Q$, would have at least twelve
antipodal real row classes. Four real sign coordinates have only
eight such classes. This is impossible. Therefore the Bloch rank is
at most two, and the preceding constructive theorem completes $V$
within $\mu_4$.
\end{proof}

The real codimension-four input is classical Hall--Verheiden
completion theory. Primary sources include E. Verheiden,
\emph{Integral and rational completions of combinatorial matrices},
J. Combinatorial Theory A 25 (1978), 267--276, and his dissertation
\url{https://thesis.caltech.edu/2899/}. The latter explicitly records
Hall's four-row theorem and its extension to seven. The criterion
and its two-dimensional quantum interpretation are presented as a
structural and algorithmic bridge; bibliographic novelty is not asserted.
The codimension-one statement is simpler: its residual has rank one
and diagonal one, and one of its columns directly supplies the
missing fourth-root column up to phase. Neither result supplies
the missing large partial frame at an unknown order.

The checker verifies the Bloch rank, phase-graph algorithm, and exact
Gaussian Gram identities at orders four, six and eight. A bounded
search of small quaternary partial frames found no rank-three input,
consistent with the classical-completion consequence. A three-axis
abstract Gaussian residual is also tested to distinguish the
residual geometry from a realizable quaternary partial frame.
